\bookchapter{Chapter 10: Mobile Payment and Life Services}

\chapterintro{The bank will send a verification code SMS to your reserved mobile number. Enter it and set a 6-digit payment password.}

\booksection{This payment password, your bank card password, and your lock screen password should ideally be three different passwords — don't set them the same for convenience.}

A little tip: if you're concerned about security, you can open a Tier II or Tier III electronic account in your bank app specifically for binding to mobile payment. Only keep your everyday spending money in this card, so even if something happens, the loss is limited.

\begin{enumerate}
  \item 2. The Payment Code Is "Whoever Sees It, Takes the Money"
  \item Many people are confused about the difference between a payment code and a collection code, and this distinction is very important.
  \item Payment Code
  \item : The barcode + QR code you show to the cashier to scan. This code is equivalent to your "electronic wallet" — whoever scans it can deduct money from your account. So
\end{enumerate}

absolutely never screenshot it and send it to anyone

\booksection{. This code automatically refreshes once per minute precisely to prevent fraudulent use.}

Collection Code

\begin{itemize}
  \item \textbf{: Others scan this code to transfer money to you. This code can be screenshotted, saved, and sent to others — it's safe.}Remember this distinction, and you can avoid many fraud traps. Anyone who asks you to "screenshot your payment code and send it over" is 100\% a scammer.\textbf{3. These Security Settings — 5 Minutes to Set Up Can Save You a Lot of Worry}Fingerprint/Face Payment
  \item \textbf{: In Alipay and WeChat's payment settings, turn on biometric payment. This way, every payment requires your fingerprint or face scan, so even if your phone is lost, the person who finds it can't spend your money.}Security Lock
\end{itemize}

: Both Alipay and WeChat have a “Security Lock” feature. Once enabled, entering the app's “Wallet” page requires an additional identity verification. Path: Alipay Me Settings Account \& Security Security Center Security Lock.

\booksection{Password-Free Payment Check}

\begin{itemize}
  \item \textbf{: Many people have signed up for various password-free payments and forgotten about them. Go to "Payment Settings" in Alipay and WeChat, find "Password-Free Payment / Auto Deduction", and see which services you've signed up for — turn off the ones you don't need.}Login Device Management
  \item \textbf{: Regularly check which devices have logged into your account. If you find an unfamiliar device, kick it out immediately. Path: Alipay Settings Account \& Security Login Device Management.}4. Bus and Subway: Your Phone Is Your Transit Card
  \item \textbf{Nowadays, very few people use physical transit cards when taking the bus or subway. There are two paths for swiping your phone at the turnstile:}Virtual Transit Card in Phone Wallet
  \item \textbf{: Open a virtual transit card directly in Apple Wallet, Huawei Wallet, or Xiaomi Wallet. The advantage of this solution is: you don't need to unlock your phone when swiping — just hold it close, the response is fastest, and it works even without a network. Some cities' virtual transit cards also offer the same transfer discounts as physical cards.}Alipay/WeChat Ride Code
\end{itemize}

\booksection{: Search for “Ride Code” in the app, and activate the bus or subway ride code for your city. The advantage of this solution is: no need to top up — it deducts directly from your balance or linked bank card. But it requires a network connection, and you need to open the app and find the code before swiping.}

I personally recommend using the virtual transit card in your phone wallet. Although you need to top up 10 or 20 yuan the first time, it's really much more convenient to use, especially when rushing for the subway — you don't have to stop to unlock your phone and find the app.

\begin{itemize}
  \item \textbf{5. Utility Bills, Food Delivery, Ride-Hailing}Water/Electricity/Gas Bills
  \item \textbf{: In Alipay or WeChat's "Utility Bills", bind your household water, electricity, and gas account numbers. Then each month when you receive a bill, just tap in to pay. The account number can usually be found on your paper bill, or you can call customer service to ask.}Food Delivery Price Comparison
\end{itemize}

: The same restaurant often has different prices and promotions on Meituan and Ele.me. Check both before placing an order — sometimes you can save several yuan. Also, Meituan's “Pin Hao Fan” feature lets you order together with friends for a noticeably cheaper price.

\booksection{Ride-Hailing Safety}

\begin{itemize}
  \item \textbf{: In Didi or Amap Ride-Hailing, remember to set up "Emergency Contact" and "Trip Sharing". When you take a ride, the app automatically shares your trip with your designated contacts. Also, for late-night rides, prioritize "Taxi" over "Express" — taxis have industry regulation and a higher safety level.}6. Handle Medical Insurance, Social Security, and Government Affairs on Your Phone
  \item \textbf{Things have changed a lot in recent years — many things that used to require queuing at a service hall can now be done in a few minutes on your phone.}Electronic Medical Insurance Certificate
  \item \textbf{: Search for "Electronic Medical Insurance Certificate" in Alipay or WeChat. After activation, you can just swipe your phone when seeing a doctor or buying medicine at a pharmacy — no need to carry your physical medical insurance card. The first time requires facial recognition verification, and afterwards a dynamic code is generated each time you use it.}Social Security Inquiry
\end{itemize}

\booksection{: Download the “Palm 12333” app or use “Citizen Center” in Alipay to check social security contribution records, pension account balances, unemployment benefit claims, and more.}

Individual Income Tax

\begin{itemize}
  \item \textbf{: The annual tax reconciliation from March to June each year can be handled on the "Individual Income Tax" app. Remember to check if there are any special additional deductions you can claim (children's education, mortgage interest, supporting the elderly, etc.) — these can all help you pay less tax or get a refund.}Remember These to Avoid Scams and Stay Safe
  \item \textbf{Police, prosecutors, banks, and customer service will never ask you to transfer money to a "safe account".}Anyone who asks you to screenshot your payment code and send it over is a scammer.
  \item \textbf{Don't click unfamiliar links, especially SMS messages like "ETC expired" or "Social security card suspended".}Online shopping refunds are returned via the original payment method — no need to provide your bank card number.
\end{itemize}

\begin{tcbwarn}{Do It Now: Give Your Payment Apps a Security Checkup}
\begin{itemize}
  \item Open Alipay and WeChat, confirm that fingerprint/face payment is enabled.
  \item Find the "Security Lock" feature and turn it on.
  \item Check the "Password-Free Payment / Auto Deduction" list, and turn off services you no longer need.
  \item View the "Login Devices" list, and remove any unfamiliar devices.
\end{itemize}

\end{tcbwarn}

\begin{tcbpractice}{Activate your electronic medical insurance certificate, so next time you visit the hospital you won't need to bring the physical card.}
\begin{enumerate}
  \item Chapter 11: Photography and Photo Management
  \item To be honest, most people's phone photos fall far short of the phone's actual potential. It's not that your skills are bad — it's that you haven't noticed many small settings and habits. This chapter skips the fluff and only covers things you can use right away.
  \item 1. Have You Made These Mistakes Too?
  \item Let's first talk about the most common problems — see how many you've hit:
  \item The lens is dirty
\end{enumerate}

\end{tcbpractice}
